% 122-36(122쪽) — 두 곡선 y=x^4-x^2, y=-x^4-x 가 둘러싸는 도형과 그 사이를 지나는 곡선 y=kx(x+1). 스캔 화질이 낮아 TikZ 로 다시 그림.
% 원문에 있는 것만: x축 눈금 -1 · 색칠한 도형 · 라벨 O, x, y, y=x^4-x^2, y=-x^4-x, y=kx(x+1).
% 가운데 곡선은 원문처럼 -1 과 0 에서 x축을 지나는 얕은 볼록 곡선으로만 그리고 눈금을 넣지 않는다
% — 10k 가 답이므로 그림에서 k 를 읽을 수 없어야 한다.
% 원문은 글자가 아주 작아 y=kx(x+1) 라벨이 색칠 영역 안에 들어가므로, 같은 배치가 되도록 그림을 크게 잡았다(가로:세로 비는 원문과 같다).
\begin{tikzpicture}[x=3.4cm, y=3.74cm, line width=0.5pt]
  % 색칠한 도형(원문) — 두 곡선 사이
  \path[fill=red!12]
    plot[domain=-1:0, samples=70] (\x, {-\x*\x*\x*\x-\x})
    -- plot[domain=0:-1, samples=70] (\x, {\x*\x*\x*\x-\x*\x}) -- cycle;
  \draw[->, >=stealth, line width=0.7pt] (-1.35,0) -- (0.42,0) node[below=1pt] {$x$};
  \draw[->, >=stealth, line width=0.7pt] (0,-0.46) -- (0,0.45) node[right=1pt] {$y$};
  % 두 곡선과 가운데 곡선
  \draw[line width=0.6pt, domain=-1.15:0.26, samples=140, smooth] plot (\x, {\x*\x*\x*\x-\x*\x});
  \draw[line width=0.6pt, domain=-1.10:0.26, samples=140, smooth] plot (\x, {-\x*\x*\x*\x-\x});
  \draw[line width=0.6pt, domain=-1:0, samples=70, smooth] plot (\x, {-0.5*\x*(\x+1)});
  \node[above=1pt, font=\small] at (-1.14,0) {$-1$};
  \node[below left=1pt and 1pt, font=\small] at (0,0) {$\pt{O}$};
  \node[anchor=south west, font=\small] at (-1.17,0.50) {$y=x^4-x^2$};
  \node[anchor=north west, font=\small] at (-1.32,-0.40) {$y=-x^4-x$};
  \node[font=\small] at (-0.573,0.21) {$y=kx(x+1)$};
\end{tikzpicture}
